\section{Geometric data and harmonic coordinates}\label{sec:data}

We begin by specifying the contraction problems under consideration.
The hypotheses distinguish the geometry needed to identify a compact
moduli space from the local estimates needed to control its metric.
They will be verified, including the entire ideal boundary, for the
eightfold constructed in Section~\ref{sec:construction}.

\subsection{Conventions}
For a projective K3 surface \(X\), write
\[
 v(F)=\ch(F)\sqrt{\operatorname{td}_X}=(r,c,s),\qquad
 ((r,c,s),(r',c',s'))=c\cdot c'-rs'-r's .
\]
Here and below \(\ch\) denotes the Chern character. We fix the
holomorphic determinant. For a smooth family of Ricci-flat Kähler forms
\(\omega_\theta\), set
\(\deg_\theta F=c_1(F)\cdot[\omega_\theta]\) and
\(\mu_\theta(F)=\deg_\theta F/\rank F\).
The harmonic metric is the real part of the \(L^2\) Hermitian pairing
on harmonic \((0,1)\)-deformations, without division by the volume of
\(X\). In particular it is half the real connection convention.

If \(\mathcal E\) is a universal topological \(K\)-class and
\(p:X\times M\to M\), our Mukai homomorphism and family indices are
\begin{equation}\label{eq:mukai-convention}
 \theta_v(z)=\bigl[p_*(\ch(\mathcal E)\sqrt{\operatorname{td}_X}\,z^\vee)\bigr]_2,
 \qquad
 I_z=p_!(z^\vee\otimes\mathcal E),\qquad
 \chi(F,z)=-(v(F),v(z)).
\end{equation}
For \(\theta_v\) the argument belongs to \(v^\perp\); the notation
\(z^\vee\) reverses its degree-two component. A class on \(X\) in a
family-index expression is pulled back without further notation.
Base line bundles may change \(\mathcal E\); they do not change
\(\theta_v|_{v^\perp}\).
The Kähler form of the actual harmonic metric has class
\begin{equation}\label{eq:harmonic-class}
 [\Omega_\theta]=2\pi^2\theta_v(b_\theta),\qquad
 b_\theta=\left(0,[\omega_\theta],
       \frac{c_1(v)\cdot[\omega_\theta]}{\rank v}\right).
\end{equation}
This is the fixed-determinant Chern--Weil formula
\cite[Proposition~4.3, (4.6)]{biswasschumacher} in the real Dolbeault
normalization. Indeed, with
\(\Delta(\mathcal E)=c_2(\mathcal E)-(r-1)c_1(\mathcal E)^2/(2r)\),
that formula reads
\(\omega_{\rm WP}/(2\pi^2)=2p_*(\Delta(\mathcal E)\omega_\theta)\).
Our form is \(\omega_{\rm WP}/2\), and expansion of
\eqref{eq:mukai-convention} gives \eqref{eq:harmonic-class}.
The computation is valid for local twisted families because their
endomorphism bundles are untwisted.

\subsection{The class of contractions}\label{subsec:hypotheses}
Fix an ample central class \(H=[\omega_0]\), a primitive positive-rank
Mukai vector \(v\), and the smooth finite-dimensional family
\(\omega_\theta\). We impose the following three hypotheses.

\medskip\noindent
\textbf{(G1) Compact chamber spaces.}
There are finitely many adjacent chambers and a connected normal
projective variety \(Z\). For all sufficiently small generic parameters,
the full moduli space \(M_\theta\) of slope-stable bundles is smooth
and projective and is also the generic Gieseker moduli space.
There is a proper surjective morphism
\(\pi_\theta:M_\theta\to Z\), constant as a complex morphism within
each chamber. Points of \(Z\) are the \(H\)-polystable classes of
vector \(v\). Every chamber object is \(H\)-semistable, and
\(\pi_\theta\) assigns its \(H\)-graded class. The \(H\)-stable locus
\(U\subset Z\) is connected and dense, and \(\pi_\theta\) is an
isomorphism over \(U\). Fix a finite Whitney stratification refining
the stable-factor decomposition types. The projective and fixed-\(H\)
gauge topologies on \(Z\) agree.

\medskip\noindent
\textbf{(G2) Absence of an ideal boundary.}
Every central ideal Hermitian--Yang--Mills object of vector \(v\) has
zero lost charge, and every \(H\)-semistable sheaf of that vector is
locally free. Thus surface gauge compactness, with the smoothly varying
background forms, has no concentration.

\medskip\noindent
\textbf{(G3) Full convergent quotient models.}
For each central class
\[
 E_p=\bigoplus_i F_i\otimes\C^{m_i},
\]
the full doubled Ext-quiver complex moment zero fibre contains a
module-simple representation of the specified dimension vector.
Here module-simple means that there is no nonzero proper submodule.
All loops, multiplicities and trivial tangent directions are retained.
The convergent symplectic coordinates give local classifying
identifications compatible with the central graded class and the
coarse analytic germ of \(Z\). In these germs,
\begin{equation}\label{eq:classifying-diagram}
 \pi_\theta f_{p,\theta}=f_{p,0}q_{p,\theta},
\end{equation}
where \(q_{p,\theta}\) is the projective quotient map to the affine
zero quotient. This is initially a statement about germs, not about
the full preimage of a target neighborhood.

\medskip
The analytic ingredients in (G3) have precise hypotheses. The
common-slope normal form
\cite[Proposition~2.1]{wuebben17} applies to locally free
\(H\)-polystable bundles on a Ricci-flat K3, with arbitrary ranks and
multiplicities. It is a convergent equivariant holomorphic change
\(\Psi(x)=x+O(\|x\|^3)\), \(D\Psi=I+O(\|x\|^2)\), identifying the
full obstruction ideal with the complex moment ideal.
The resulting harmonic comparison
\cite[Corollary~2.2]{wuebben17}, obtained from the local harmonic
reduction \cite[Theorem~1.1 and Sections~3,5]{wuebben15} and
finite balancing \cite[Theorem~2.1]{wuebben14}, gives
\begin{equation}\label{eq:relative}
 (1-e_p)g_{p,\zeta_p(\theta)}
 \ \leq\ f_{p,\theta}^*g_\theta\
 \leq (1+e_p)g_{p,\zeta_p(\theta)},\qquad
 e_p=C_p(|\theta|+\|x\|^2).
\end{equation}
The inequalities concern tangents to actual smooth quotient pieces.
The local theorem also covers all sufficiently nearby actual solutions.

The finite quotient metric in \eqref{eq:relative} is induced by the
central harmonic norm on \(\Ext^1(E_p,E_p)\). Its real moment covector,
paired by ordinary multiplicity trace, is
\begin{equation}\label{eq:levels}
 \zeta_{p,i}(\theta)=2\pi\bigl(r_i\mu_\theta(v)-\deg_\theta F_i\bigr),
 \qquad \sum_i m_i\zeta_{p,i}=0,\qquad r_i=\rank F_i.
\end{equation}
The compact group is \(\prod_iU(m_i)\), modulo the ineffective common
scalar. The rank weights instead occur in the spectral mass
\(\Vol(\omega_0)\sum_i r_i\tr(\xi_i^*\xi_i)\), whose small quadratic
form is \(\|[\xi,x]\|^2\). These are different pairings.
Constants in \eqref{eq:relative} depend on the fixed background and
its complementary elliptic gaps. They do not depend on ratios of
positive finite eigenvalues or on a nearby stabilizer. Formal
quadraticity alone would not give this estimate; the controlled
convergent coordinates are essential.

\subsection{What the reconstruction retains}\label{subsec:retained}
Retain \(U\) as a geometric space with its interpretation by stable
bundles and local classifying families. Retain also the quotient
domains with their actual families in this same moduli problem.
On two domains, their overlap is the locus of pairs representing
isomorphic bundles. Injective coarse classifying maps identify its
two projections with open subsets and give a unique transition map.
We take cochains only after these overlap spaces and maps have
been specified. In particular a bare cochain complex of \(U\) is not
enough to define the overlaps.

These geometric maps retain stratum monodromy and all multiple
compatibilities. Local fibre types alone do not supply them.
Specialization of exceptional classes will be computed from a
separate proper family in Section~\ref{subsec:specialization}.
For the smaller middle-degree calculation, the global datum is
instead the marked ring \(H^*(M\setminus E,\Z)\), where \(E\) is the
chosen compact exceptional core. This complement generally contains
exceptional fibres over other strata. In our eightfold it is larger
than \(\pi^{-1}(U)\).
